\section{Introduction}

Deep learning models for skin lesion classification have achieved dermatologist-level accuracy on curated benchmarks, yet their performance is not uniform across skin tones. Systematic evaluations have revealed that diagnostic models trained predominantly on lighter skin tones exhibit degraded sensitivity and specificity on darker skin presentations \cite{Esteva2017, Daneshjou2022}. Groh et al.~\cite{Groh2021} documented substantial variation in model performance across Fitzpatrick skin types, underscoring that reported aggregate metrics can mask critical disparities. These findings raise concerns about the equitable clinical deployment of dermatological AI, particularly for populations that already face barriers to dermatologic care \cite{Adamson2018}.

The Dermatology Data Initiative (DDI) dataset represents a recent effort to address this gap by providing annotated dermatoscopic images \cite{Tan2019} with Fitzpatrick skin type labels \cite{Fitzpatrick1988}. However, the dataset remains small relative to large-scale benchmarks, and dark-skin melanoma cases are particularly scarce. The HAM10000 dataset \cite{Tschandl2018}, while substantially larger, lacks skin type annotations and skews toward lighter skin presentations. This class imbalance means that models pretrained on such data may encode skin-tone priors that propagate into downstream fine-tuning, especially when dark-skin melanoma examples number fewer than 50.

Synthetic data augmentation offers a practical approach to mitigating training-set imbalance without requiring additional clinical data collection. Modern augmentation libraries such as Albumentations \cite{Buslaev2020} support principled color-space and geometric perturbations that can expand the effective diversity of underrepresented subgroups. Fairness-aware augmentation strategies have been explored in general machine learning contexts \cite{Mehrabi2021}, though their application to dermatological classification remains under-studied.

In this work, we investigate synthetic augmentation as a fairness intervention for melanoma detection across skin tones on the DDI dataset. Our contributions are threefold: (1) we quantify the skin-tone bias gap in an EfficientNet-B0 model pretrained on HAM10000 and fine-tuned on DDI, revealing near-chance performance on dark-skin melanoma; (2) we evaluate a leakage-free, train-only synthetic augmentation pipeline ($n\,{=}\,29$ source dark melanomas, 290 synthetics) and show simple photometric transforms do not significantly improve Dark AUROC ($0.472 \to 0.463$, $p\,{=}\,0.524$), demonstrating limits of color-jitter fairness interventions; and (3) we conduct an ablation study over augmentation factors (1$\times$--10$\times$) to identify near-optimal augmentation strength. Main results are summarized in Table~\ref{tab:main_results}, with detailed subgroup AUROC values in Table~\ref{tab:subgroup_auroc} and visualized in Figure~\ref{fig:subgroup_auroc}.
